\pdfoutput=1
\documentclass[11pt]{amsart}
\usepackage[T1]{fontenc}
\usepackage[utf8]{inputenc}
\usepackage{lmodern}
\usepackage[margin=1.15in]{geometry}
\usepackage{amsmath,amssymb,amsthm,mathtools}
\usepackage{booktabs}
\usepackage[section]{placeins}
\usepackage{microtype}
\usepackage{enumitem}
\usepackage[hidelinks]{hyperref}
\usepackage{xurl}
\usepackage{orcidlink}
\usepackage[capitalise,nameinlink]{cleveref}

\newtheorem{theorem}{Theorem}[section]
\newtheorem{lemma}[theorem]{Lemma}
\newtheorem{corollary}[theorem]{Corollary}
\newtheorem{proposition}[theorem]{Proposition}
\theoremstyle{remark}
\newtheorem{remark}[theorem]{Remark}

\newcommand{\leanus}{\textunderscore\allowbreak}
\newcommand{\lean}[1]{{\ttfamily\let\_\leanus #1}}
\newcommand{\code}[1]{\texttt{#1}}
\newcommand{\fd}{f_{\delta\ge 3}}
\newcommand{\fc}{f_{\mathrm{cubic}}}
\newcommand{\Hseven}{H_7}
\newcommand{\Hfifteen}{H_{15}}
\newcommand{\Height}{H_8^{*}}
\newcommand{\TC}{\mathrm{TC}}

\title[Tutte--Coxeter lifts without 32-cycles]{The smallest Tutte--Coxeter lifts\\ without cycles of length 32}

\author[A. Marszałek]{Andrzej Marszałek\,\orcidlink{0009-0006-6148-2965}}
\address{Independent Researcher}
\urladdr{https://github.com/hipotures/conjectures}

\subjclass[2020]{05C38, 05C35, 68V15}
\keywords{Erdős–Gyárfás conjecture, power-of-two cycles, Tutte–Coxeter graph, vertex replacement, SAT, DRAT proofs}

\date{28 November 2026}

\begin{document}

\begin{abstract}
The best known graphs of minimum degree three without cycles of length $4$,
$8$, $16$ and $32$ are obtained from the Tutte--Coxeter graph by replacing
every vertex with one of three small gadgets: Henry's cubic graph on $370$
vertices and our earlier graph on $345$ vertices. We show that neither
construction can be improved within its family. Every such lift without a
cycle of length $32$ has at least $345$ vertices, and at least $370$ if it is
cubic. The proof first shows that, for these gadgets, a lift has a $32$-cycle
exactly when some cycle of the base is too light. A double-counting identity
over the $90$ eight-cycles of the Tutte--Coxeter graph, proved in Lean~4, then
leaves three cases. Each of them is a SAT instance that CaDiCaL refutes; every
refutation is a DRAT proof checked by \code{drat-trim}. An independent test
suite checks the encoding against lifts built as graphs.
\end{abstract}

\maketitle

%======================================================================
\section{Introduction}

\begin{quote}
\textbf{Conjecture} (Erdős and Gyárfás, 1995; see~\cite{Erdos1997}).
Every finite simple graph of minimum degree at least three contains a cycle
whose length is a power of two.
\end{quote}

For $k\ge 2$ let $\fd(k)$ and $\fc(k)$ be the least order of a graph of
minimum degree at least $3$, respectively of a cubic graph, with no cycle of
length $4,8,\dots,2^k$~\cite{Garcia2026,Henry2026,Marszalek2026a}. For $k=5$
the best bounds come from one construction. Take the Tutte--Coxeter graph
$\TC$, the $(3,8)$-cage on $30$ vertices~\cite{Tutte1947}, and replace each
vertex by a small three-port gadget. Henry~\cite{Henry2026} obtained
$\fc(5)\le 370$ with a construction of this type. Our earlier
paper~\cite{Marszalek2026a} gives a $370$-vertex lift with the gadgets
$\Hseven$ and $\Hfifteen$ of García~\cite{Garcia2026}, and a $345$-vertex lift
with $\Hfifteen$ and a new gadget $\Height$ on eight vertices, which shows
$\fd(5)\le 345$. Both lifts have no cycle of length $4$, $8$, $16$ or $32$.

Can a cleverer arrangement of the same gadgets on the same base do better?
There are $18^{30}$ ways to place a gadget and a port bijection at each vertex,
so the question is not settled by trying them.

\begin{theorem}\label{thm:main}
Let $L$ be a lift of the Tutte--Coxeter graph in which every vertex is
replaced by $\Hseven$, $\Height$ or $\Hfifteen$, with arbitrary port
bijections. If $L$ has no cycle of length $32$, then
\begin{enumerate}[label=\textup{(\roman*)}]
\item $|V(L)|\ge 345$, and
\item $|V(L)|\ge 370$ if $L$ is cubic, that is, if no gadget is $\Height$.
\end{enumerate}
Both bounds are attained by lifts with no cycle of length $4$, $8$, $16$ or
$32$~\cite{Marszalek2026a}.
\end{theorem}

So the constructions are optimal in their family, and an improvement of
$\fd(5)\le 345$ or $\fc(5)\le 370$ needs a different base or different gadgets.
The theorem says nothing about such graphs, and nothing about lower bounds
for $\fd(5)$ or $\fc(5)$ themselves.

Forbidding only $32$-cycles costs nothing here: for these gadgets a lift
without $32$-cycles automatically has no cycle of length at most $32$ that
meets two gadgets (\cref{lem:c32}), so it also has no cycle of length $4$, $8$
or $16$.

\subsection*{Nature of the proof}
The proof has three parts. \Cref{sec:lifts} reduces the absence of
$32$-cycles to a condition on the $90$ eight-cycles of $\TC$; it uses a short
table of path lengths inside the gadgets, computed by exhaustive enumeration.
\Cref{sec:counting} proves a double-counting identity over these eight-cycles.
It is formalized in Lean~4~\cite{Lean4} in the library of our earlier
paper~\cite{Marszalek2026a}, and it leaves three cases. \Cref{sec:sat} encodes
each case as a CNF formula. CaDiCaL~\cite{CaDiCaL2024} shows that all three are
unsatisfiable and produces DRAT proofs, and \code{drat-trim}~\cite{DRATtrim2014}
checks every proof against the formula. What is not machine-checked is that the
CNF formulas say what \cref{sec:sat} says they say. \Cref{sec:validation}
describes the tests that address this.

%======================================================================
\section{Lifts and 32-cycles}\label{sec:lifts}

We use the setting of~\cite[Section~2]{Marszalek2026a}. A
\emph{three-port gadget} is a graph $H$ with three ports $p_0,p_1,p_2$ of
degree $2$ and all other vertices of degree at least $3$. Given a cubic graph
$B$, a gadget $H_v$ at every vertex $v$ and a bijection $\sigma_v$ from the
three edges at $v$ to $\{0,1,2\}$, the \emph{lift} $B[H,\sigma]$ is the
disjoint union of the $H_v$ together with an edge from port $\sigma_u(uv)$ of
$H_u$ to port $\sigma_v(uv)$ of $H_v$ for every edge $uv$ of $B$. For a cycle
$D=(v_1,\dots,v_\ell)$ of $B$ put $a_i=\sigma_{v_i}(v_{i-1}v_i)$,
$b_i=\sigma_{v_i}(v_iv_{i+1})$ and
\[
  w(D)=\ell+\sum_{i=1}^{\ell} d_{H_{v_i}}(a_i,b_i).
\]

\begin{table}[ht]
\centering
\caption{Lengths of paths between two ports of each gadget. Every set is an
interval. $d(0,1)=d(0,2)$ in all three gadgets.}
\label{tab:paths}
\begin{tabular}{@{}lrrrr@{}}
\toprule
Gadget & $|V|$ & ports $0,1$ and $0,2$ & ports $1,2$ & $s(H)=d(1,2)+2d(0,1)$\\
\midrule
$\Hseven$   & 7  & $2,\dots,6$  & $3,\dots,6$  & 7\\
$\Height$   & 8  & $3,\dots,7$  & $2,\dots,7$  & 8\\
$\Hfifteen$ & 15 & $3,\dots,14$ & $5,\dots,14$ & 11\\
\bottomrule
\end{tabular}
\end{table}

The gadgets are those of~\cite[Section~3]{Marszalek2026a}; their edge lists
are in \code{data/gadgets.json}. \Cref{tab:paths} was obtained by listing all
port-to-port paths.

\begin{lemma}\label{lem:c32}
Let $B$ be a cubic graph of girth at least $8$ and $L$ a lift of $B$ by the
gadgets $\Hseven$, $\Height$, $\Hfifteen$. Then $L$ has a cycle of length $32$
if and only if $w(D)\le 32$ for some cycle $D$ of $B$. If $w(D)\ge 33$ for
every cycle $D$ of $B$, then $L$ has no cycle of length $4$, $8$, $16$ or $32$.
\end{lemma}

\begin{proof}
A cycle inside one gadget has at most $15$ vertices, and the gadgets have no
cycle of length $4$ or $8$~\cite[Table~2]{Marszalek2026a}. A cycle $C$ of $L$
that meets two gadgets visits each gadget in at most one run, and contracting
the gadgets turns $C$ into a cycle $D=(v_1,\dots,v_\ell)$ of
$B$~\cite[Lemma~2.1]{Marszalek2026a}. The run in $H_{v_i}$ is a path from port
$a_i$ to port $b_i$, so $|C|=\ell+\sum_i |P_i|$ with $P_i$ that path.
Conversely, any choice of $a_i$--$b_i$ paths $P_i$ in the gadgets $H_{v_i}$,
joined by the $\ell$ edges between consecutive gadgets, is a cycle of $L$,
because the gadgets are disjoint. By \cref{tab:paths} the possible lengths of
$P_i$ form an interval $[d_i,M_i]$ with $d_i=d_{H_{v_i}}(a_i,b_i)$ and
$M_i\ge 6$. A sum of integer intervals is an interval, so the lengths of the
cycles of $L$ that contract to $D$ are exactly the integers from $w(D)$ to
$\ell+\sum_i M_i\ge 7\ell\ge 56$. Hence $32$ is such a length if and only if
$w(D)\le 32$. Both statements follow.
\end{proof}

At a vertex $v$ of a cycle $D$, the edge of $v$ that is not on $D$ goes to some
port $o$; the cycle passes through the other two. If $o=0$, the contribution
of $v$ to $w(D)$ is $d(1,2)$, and otherwise it is $d(0,1)$, since
$d(0,1)=d(0,2)$. So $w(D)$ depends only on the gadgets and on the neighbour
$q(v)$ joined to port $0$:
\begin{center}
\begin{tabular}{@{}lccc@{}}
\toprule
& $\Hseven$ & $\Height$ & $\Hfifteen$\\
\midrule
edge to $q(v)$ off $D$ & 3 & 2 & 5\\
edge to $q(v)$ on $D$  & 2 & 3 & 3\\
\bottomrule
\end{tabular}
\end{center}
Every contribution is at least $2$, so $w(D)\ge 3\ell\ge 36$ when $\ell\ge 12$.
The Tutte--Coxeter graph is bipartite of girth $8$, so only its $90$ cycles of
length $8$ and $72$ cycles of length $10$ can be light. For the lower bound we
need only the eight-cycles.

\begin{corollary}\label{cor:necessary}
If a lift of $\TC$ by these gadgets has no $32$-cycle, then $w(D)\ge 33$ for
each of the $90$ eight-cycles $D$ of $\TC$.
\end{corollary}

%======================================================================
\section{Counting}\label{sec:counting}

\begin{lemma}\label{lem:incidence}
For every vertex $v$ of $\TC$ and every edge $e$ at $v$, exactly $8$ of the
$90$ eight-cycles pass through $v$ and avoid $e$.
\end{lemma}

\begin{proof}
The automorphism group of $\TC$, of order $1440$, acts transitively on pairs
$(v,e)$ with $e$ an edge at $v$, since $\TC$ is arc-transitive (indeed
$5$-arc-transitive~\cite{Tutte1947}). There are $90$ such pairs. Each
eight-cycle passes through $8$ vertices and avoids exactly one edge at each of
them, giving $8$ pairs per cycle and $720$ incidences in total. By
transitivity each pair receives $720/90=8$.
\end{proof}

Let $a$, $b$ and $c$ be the numbers of copies of $\Hfifteen$, $\Height$ and
$\Hseven$ in a lift, so $a+b+c=30$ and the lift has $15a+8b+7c=240+7a-c$
vertices.

\begin{lemma}\label{lem:identity}
For every lift of $\TC$ by these gadgets,
\[
  \sum_{D}\bigl(w(D)-33\bigr) \;=\; 88a+64b+56c-2250 \;=\; 24a-8c-330,
\]
where $D$ runs over the $90$ eight-cycles of $\TC$.
\end{lemma}

\begin{proof}
The sum of $w(D)$ is $90\cdot 8$ plus the sum, over pairs $(v,e)$, of the
contribution of $v$ when $e$ is the edge of $v$ off the cycle, counted $8$
times by \cref{lem:incidence}. Summed over the three edges at $v$, the
contributions give $s(H_v)$ of \cref{tab:paths}. Hence the sum of $w(D)$ is
$720+8(11a+8b+7c)$, and subtracting $90\cdot 33$ gives the identity. The second
form uses $b=30-a-c$.
\end{proof}

\begin{corollary}\label{cor:counting}
If a lift of $\TC$ by these gadgets has no $32$-cycle, then
\begin{enumerate}[label=\textup{(\alph*)}]
\item at most $24a-8c-330$ of the eight-cycles have $w(D)\ge 34$;
\item $3a-c\ge 42$, and the lift has at least $282+4a$ vertices.
\end{enumerate}
\end{corollary}

\begin{proof}
By \cref{cor:necessary} every term of \cref{lem:identity} is non-negative,
and each eight-cycle with $w(D)\ge 34$ contributes at least $1$. This gives (a)
and $24a-8c\ge 330$. The left side is divisible by $8$, so $24a-8c\ge 336$,
that is, $3a-c\ge 42$. Then $240+7a-c=240+(3a-c)+4a\ge 282+4a$.
\end{proof}

The formal versions are \lean{tc\_incidence\_eq\_8},
\lean{tc\_double\_counting\_eq} and \lean{tc\_slack\_card\_sound} in
\code{lean/ErdosGyarfas/TutteCoxeter.lean}. They state the same facts for an
explicit list of $90$ distinct eight-cycles, with the contributions of the
table above, and \lean{tcVertexWeight\_eq\_liftWeight} connects these
contributions to the port distances of the gadgets used in the lift
construction of~\cite{Marszalek2026a}. The test suite checks that the list in
Lean is the set of eight-cycles used by the encoder.

\begin{proposition}\label{prop:sat}
No lift of $\TC$ satisfies the conclusions of \cref{cor:necessary} and
\cref{cor:counting}(a) in any of the cases
\begin{itemize}[leftmargin=2.2em]
\item[\textsf{A}:] gadgets $\Hfifteen,\Height$, at most $14$ copies of $\Hfifteen$;
\item[\textsf{B}:] all three gadgets, at most $15$ copies of $\Hfifteen$, at least one $\Hseven$;
\item[\textsf{C}:] gadgets $\Hfifteen,\Hseven$, at most $19$ copies of $\Hfifteen$.
\end{itemize}
\end{proposition}

In each case the bound of \cref{cor:counting}(a), maximized over the allowed
counts, is $6$, $22$ and $38$ respectively. \Cref{sec:sat,sec:results}
prove the proposition.

\begin{proof}[Proof of \cref{thm:main}, given \cref{prop:sat}]
Suppose a lift without $32$-cycles has $n$ vertices.

(i) Let $n\le 344$. By \cref{cor:counting}(b), $a\le 15$. If $a\le 14$, then
$c\le 3a-42\le 0$, so $c=0$ and the lift uses only $\Hfifteen$ and $\Height$,
with at most $14$ copies of $\Hfifteen$: this is case~\textsf{A}. If $a=15$,
then $n=345-c\le 344$ gives $c\ge 1$: this is case~\textsf{B}.

(ii) If the lift is cubic, then $b=0$ and $c=30-a$, so (b) gives $a\ge 18$ and
$n=210+8a$. If $n<370$ then $a\le 19$: this is case~\textsf{C}.

By \cref{prop:sat} none of the three cases occurs.
\end{proof}

%======================================================================
\section{The SAT encoding}\label{sec:sat}

Number the vertices of $\TC$ by its LCF code $[-13,-9,7,-7,9,13]^5$, and let
$(n_{v,0},n_{v,1},n_{v,2})$ be the neighbours of $v$ in increasing order.
For each vertex $v$, allowed gadget $g$ and $p\in\{0,1,2\}$ a variable
$x_{v,g,p}$ says that $H_v=g$ and $q(v)=n_{v,p}$. The formula contains:
\begin{enumerate}[label=(\arabic*)]
\item exactly one $x_{v,g,p}$ for every $v$ (pairwise exclusions);
\item for each eight-cycle $D$, the constraint $w(D)\ge 33$. Write the
  contribution of each vertex of $D$ as $2+e$ with $e\in\{0,1,3\}$. A variable
  for each value of $e$ at each pair $(v,\text{edge off }D)$ is defined as the
  disjunction of the states $x_{v,g,p}$ that give it. Prefix variables
  $s_{i,t}$, meaning that the first $i$ vertices of $D$ contribute $e$-sum at
  least $t$, are defined by full equivalences in a sequential layout, and the
  formula asserts $s_{8,9}$, which is $w(D)\ge 33$;
\item a sequential-counter encoding~\cite{Sinz2005} of the bounds on the
  number of copies of $\Hfifteen$ and $\Hseven$;
\item the bound of \cref{cor:counting}(a): with $s_{8,10}$ (that is,
  $w(D)\ge 34$) for each eight-cycle, at most $S$ of these $90$ variables are
  true, where $S=6,22,38$ as above;
\item symmetry breaking: in case \textsf{A} vertex $0$ is $\Height$, in cases
  \textsf{B} and \textsf{C} it is $\Hseven$, and in all cases $q(0)=1$.
\end{enumerate}
The symmetry breaking loses nothing. Each case contains at least one vertex
with the named gadget, and an automorphism of $\TC$ maps that vertex and its
neighbour $q(v)$ to the arc $(0,1)$, transforming the lift into an isomorphic
one with the same counts and the same weights. Case~\textsf{C} is split
further into $56$ disjoint subcases (cube and
conquer~\cite{HeuleKullmannWieringaBiere2012}): by the states of the
neighbours $1$ and $29$ of vertex $0$ ($36$ pairs), and for the four pairs in
which both are $\Hfifteen$ with $q(1)\neq 0$ and $q(29)\neq 0$, also by the
state of the third neighbour $17$. A test checks that the $56$ subcases cover each of the
$216$ states of $(1,29,17)$ exactly once.

%======================================================================
\section{Results}\label{sec:results}

\Cref{tab:results} lists the instances. All $58$ unsatisfiable formulas
were refuted by CaDiCaL, and \code{drat-trim} reported \code{s VERIFIED} for
each of the $58$ proofs. The longest solver run took $1298$ seconds and the
longest check $1195$ seconds. The two controls are satisfiable, and
their models are the optimal lifts described in \cref{sec:validation}. The
whole run took $50$ minutes on $16$ cores.

\begin{table}[ht]
\centering
\caption{The SAT instances. $S$ is the bound of \cref{cor:counting}(a). Times
are the running times of the single-threaded CaDiCaL runs and of
\code{drat-trim} on their proofs, summed over the files of a case; variables
and clauses are per file. Machine: AMD Ryzen 9 7950X3D 16-Core Processor, 16 threads, 16 parallel runs, 50 minutes of wall time.}
\label{tab:results}
\small
\begin{tabular}{@{}lrrrrrrl@{}}
\toprule
case & files & variables & clauses & $S$ & solve (s) & check (s) & result\\
\midrule
\textsf{A} & 1 & 7{,}570 & 35{,}999 & 6 & 30 & 32 & UNSAT, all checked\\
\textsf{B} & 1 & 9{,}984 & 41{,}590 & 22 & 1{,}239 & 1{,}195 & UNSAT, all checked\\
\textsf{C} & 56 & 10{,}563 & 41{,}950 & 38 & 14{,}261 & 15{,}087 & UNSAT, all checked\\
control, \textsf{A} with $a\le 15$ & 1 & 9{,}735 & 40{,}304 & 30 & 2 & -- & SAT\\
control, \textsf{C} with $a\le 20$ & 1 & 13{,}440 & 47{,}669 & 70 & 22 & -- & SAT\\
\bottomrule
\end{tabular}
\end{table}


%======================================================================
\section{Validation}\label{sec:validation}

The script \code{test\_encoding.py} uses only the Python standard library and
no SAT solver; it runs in about $15$ seconds. It does not reuse the encoder's
model of the problem:
\begin{enumerate}[label=(\alph*)]
\item It rebuilds $\TC$ as the incidence graph of the generalized quadrangle
  $\mathrm{GQ}(2,2)$ and finds an isomorphism to the encoder's graph. It
  recounts the cycles of length $8$, $10$ and $12$ ($90$, $72$ and $300$),
  checks \cref{lem:incidence}, finds all $1440$ automorphisms, and checks that
  the eight-cycles agree with the Lean list.
\item It recomputes \cref{tab:paths} and the cycle lengths of the gadgets.
\item It builds $300$ lifts near the optimum as graphs, starting from the
  $345$- and $370$-vertex lifts of~\cite{Marszalek2026a}, applying random
  automorphisms and changing a few gadgets and port bijections, including the
  choice between ports $1$ and $2$. For each lift it computes the shortest
  cycle through an edge between two gadgets by breadth-first search. It checks
  that this equals the smallest $w(D)$ over the $8$- and $10$-cycles, and that
  a small DPLL solver finds the formula made of parts~(1) and~(2) satisfiable under the
  lift's states exactly when all eight-cycles have $w(D)\ge 33$ ($48$ lifts
  accepted, $252$ rejected). For every rejected lift it constructs an explicit
  $32$-cycle, as in the proof of \cref{lem:c32}, and checks it edge by edge.
  Every accepted lift also satisfies the full formula with its own counts,
  including constraint~(4), and, after an automorphism moves one of its
  vertices to $0$, constraint~(5).
\item It checks the cardinality encoding exhaustively for up to six variables.
\end{enumerate}
We also changed the encoder on purpose: a bound of $32$ instead of $33$,
swapped port distances, and a shifted edge index. Each change made the test
fail. Two satisfiable controls give the optimal lifts: \textsf{A} with at most
$15$ copies of $\Hfifteen$, and \textsf{C} with at most $20$. The script
\code{prove.py} decodes their models and checks them as graphs: they have
$345$ and $370$ vertices, and their shortest cycle between two gadgets has
length $33$.

\begin{remark}
The $10$-cycles are not needed. In a search over lifts with all eight-cycle
weights at least $33$ we found no $10$-cycle of weight below $35$, and
encoding only the eight-cycles keeps the formula a relaxation: unsatisfiability
of the relaxation rules out every lift.
\end{remark}

%======================================================================
\section{Reproducibility}\label{sec:repro}

Everything is in the repository \url{https://github.com/hipotures/conjectures},
in the folder \nolinkurl{results/tutte-coxeter-lifts-without-32-cycles}, in the
public domain (CC0); this paper describes release \code{v1.2.0}, archived on
Zenodo as \url{https://doi.org/10.5281/zenodo.24624865}.
\code{python3 test\_encoding.py} runs the tests. \code{python3 prove.py}
downloads CaDiCaL~2.1.3 and \code{drat-trim} (commit \code{2e3b2dc}), checks
their SHA-256 hashes and builds them, writes all $60$ CNF files, solves them,
checks every DRAT proof, checks the two controls, and writes
\code{data/results.json} with the SHA-256 hash of every CNF file. The Lean
library is built with \code{lake build} in the folder \code{lean} at the root
of the repository, which is shared with~\cite{Marszalek2026a}.

\section*{Use of AI tools}
This work was carried out with substantial assistance from AI systems, mainly
Gemini~4~Neon (Google) and Claude Opus~5.5 (Anthropic), used through the
Antigravity and Claude Code environments. They were used to write and run the
encoder, the tests and the Lean proofs, and to draft the text. AI systems are
not authors of this paper. The author takes responsibility for the content and
for the attribution of prior work.

\begin{thebibliography}{99}

\bibitem{CaDiCaL2024}
A.~Biere, T.~Faller, K.~Fazekas, M.~Fleury, N.~Froleyks and F.~Pollitt,
\emph{CaDiCaL 2.0}, in: Computer Aided Verification -- CAV 2024, Lecture Notes
in Comput. Sci. \textbf{14681}, Springer, 2024, 133--152,
\url{https://doi.org/10.1007/978-3-031-65627-9_7}.

\bibitem{Erdos1997}
P.~Erdős, \emph{Some old and new problems in various branches of combinatorics},
Discrete Math. \textbf{165/166} (1997), 227--231.

\bibitem{Garcia2026}
D.~García, \emph{Small graphs without power-of-two cycles: a lower bound of 24,
a correction to a construction of Exoo, and explicit bounds for $f(k)$},
arXiv:2609.04686, 2026.

\bibitem{Henry2026}
K.~Henry, \emph{Improved cubic graph constructions avoiding power-of-two cycle
lengths}, preprint, Zenodo, 2026, \url{https://doi.org/10.5281/zenodo.23886386}.

\bibitem{HeuleKullmannWieringaBiere2012}
M.~J.~H.~Heule, O.~Kullmann, S.~Wieringa and A.~Biere, \emph{Cube and conquer:
guiding CDCL SAT solvers by lookaheads}, in: Hardware and Software:
Verification and Testing -- HVC 2011, Lecture Notes in Comput. Sci.
\textbf{7261}, Springer, 2012, 50--65,
\url{https://doi.org/10.1007/978-3-642-34188-5_8}.

\bibitem{Marszalek2026a}
A.~Marszałek, \emph{Small graphs without cycles of length a power of two: new
bounds with kernel-checked proofs}, preprint, 2026,
\url{https://doi.org/10.5281/zenodo.24588313}.

\bibitem{Lean4}
L.~de Moura and S.~Ullrich, \emph{The Lean 4 theorem prover and programming
language}, in: Automated Deduction -- CADE 28, Lecture Notes in Comput. Sci.
\textbf{12699}, Springer, 2021, 625--635.

\bibitem{Sinz2005}
C.~Sinz, \emph{Towards an optimal CNF encoding of Boolean cardinality
constraints}, in: Principles and Practice of Constraint Programming -- CP 2005,
Lecture Notes in Comput. Sci. \textbf{3709}, Springer, 2005, 827--831,
\url{https://doi.org/10.1007/11564751_73}.

\bibitem{Tutte1947}
W.~T.~Tutte, \emph{A family of cubical graphs}, Proc. Cambridge Philos. Soc.
\textbf{43} (1947), 459--474.

\bibitem{DRATtrim2014}
N.~Wetzler, M.~J.~H.~Heule and W.~A.~Hunt~Jr., \emph{DRAT-trim: efficient
checking and trimming using expressive clausal proofs}, in: Theory and
Applications of Satisfiability Testing -- SAT 2014, Lecture Notes in Comput.
Sci. \textbf{8561}, Springer, 2014, 422--429,
\url{https://doi.org/10.1007/978-3-319-09284-3_31}.

\end{thebibliography}

\end{document}
